% Pursuing Stacks (cotes 134-2 to 134-8): archive cote 134-4, site page 59 (PDF page 60), typescript page 147.
% typescript. Draft compared: ps3.tex:3009-3066.
% AI-transcribed by Claude Opus 5.5 (claude-opus-5-5, adaptive thinking, effort high), Message Batches API, batch ps2 (msgbatch_01NVwr1Yuo8qJf5rykeGBXPK), 2026-11-12; not yet checked by eye.
\scanpage{134-4/59}
\mspage{147}
matter -  the way topological spaces are used as ``homotopy models'' in
standart homotopy theory,  it is \uline{pathwise} \uline{connected} components that count,
and not the topological ones.  In terms of these,  there is still a canoni%
% NOTE: the right edge of the scan cuts the end of this line; a hyphen after "canoni" may be hidden
cal functor
\[ \hand{\pi}_0 : \quad (Spaces) \longrightarrow (Sets) \quad , \]
this functor however is no longer left adjoint of the functor in opposite
direction\hand{,} associating to every set E,  the corresponding discrete topolo-%
gical space. (To get an adjunction formula,  we should have to restrict
to the category of \inserted[typed]{pathwise} locally connected topological spaces.)  It does'nt
matter  visibly  -  all that's being used is that \hand{\pi}_0 commutes to arbitra-%
ry sums,  and takes pathwise 0-connected spaces into one-point sets.

This suggests that we should generalize the notions around the ``canon\hand{i}%
cal'' homotopy structure on a category M,  to the case of a category which
need not satisfy the exacting conditions of total 0-connectedness,  by
introducing as an \uline{extra\xout{\unsure{s}}}  \uline{structure}  upon M  a given functor
\[ \hand{\pi}_0 : M \rightarrow (Sets) \quad , \]
subject \inserted{possibly} to suitable restrictions.  The first which comes to mind here is
commutation wi th sums  -  it does'nt seem though we've had to use this
property so far.  All we've used occasionally was existence of finite pro-%
ducts in M,  and commutation of \hand{\pi}_0 to these.

If we think of M as a would-be modelizer,  and therefore endowed with
a hoped-for functor
\[ M \longrightarrow (Hot) \quad , \]
there is a natural functor \hand{\pi}_0 indeed on M ,  namely the composition
\[ M \rightarrow (Hot) \longrightarrow (Sets) \quad , \]
where the canonical functor
\[ \hand{\pi}_0 :\xout{x}(Hot) \longrightarrow (Sets) \]
is deduced from the \hand{\pi}_0-functor (Cat) $\longrightarrow$ (Sets)  considered previously,
by factorization through the localized category (Hot) of (Cat).  Thus``the
least we would expect'' from a category M for being eligible as a modelizer\unsure{,}
is that there should be a natural functor \hand{\pi}_0 around,  corresponding to the
intuition of connected components.  In case of a``canonical''modelizer  \hand{M}
(maybe we should say rather:  canonical wi th respect to a given \hand{\pi}_0 ) ,
there is the feeling that the functor  M $\rightarrow$ (Hot)  we are after could even-%
tually be squeezed out from just \hand{\pi}_0 ,  and that it could be viewed as some%
% NOTE: the right edge of the scan cuts the end of this line; a hyphen after "some" may be hidden
thing like a ``total left derived functor'' of the functor \hand{\pi}_0 .  But this
for the time being is still thin air ...

\margin{59}
